% Nag is used for debug
\RequirePackage[l2tabu, orthodox]{nag}
\documentclass[a4paper,french,twoside,10pt]{report}
%\documentclass[a4paper,twoside,10pt]{article}
%\usepackage[bookman]{quotchap}

\usepackage[bindingoffset=1.5cm,centering,includeheadfoot,margin=3cm]{geometry}

%\frenchspacing
\usepackage[frenchb]{babel}
\usepackage[T1]{fontenc}
\usepackage[utf8]{inputenc}
\usepackage[pdftex]{graphicx}
\usepackage{subfig}

\usepackage[protrusion=true,expansion=true]{microtype}

\usepackage{fancybox}
\usepackage{shadow}
\usepackage{xspace}
\usepackage{acronym}

% Maths enhancing
\usepackage{amsmath}
\usepackage{amssymb}
\usepackage{amsthm}
\usepackage{dsfont}

\usepackage{algorithm2e}

% JP : automatic line return after paragraph
\makeatletter
\renewcommand\paragraph{\@startsection{paragraph}{4}{\z@}%
    {-3.25ex\@plus -1ex \@minus -.2ex}%
    {1.5ex \@plus .2ex}%
{\normalfont\normalsize\bfseries}}
\makeatother

% JP : use amsthm package and define some useful commands
\newtheorem{theorem}{Theorem}
\newtheorem{definition}{Definition}
\newtheorem{proposition}{Proposition}

%\usepackage[right]{eurosym}
\usepackage{lmodern} % new latin extended computer modern font}, use with T1
%\usepackage{pxfonts} % has replaced palatino
%\usepackage{cmbright}

\usepackage{listings}

%\usepackage{booktabs}
%\usepackage{array}

\usepackage[allowlitunits]{siunitx}
% Try $20\micro\seconds$

%\addtolength{\oddsidemargin}{-1cm}
%\addtolength{\evensidemargin}{-1cm}
%\addtolength{\textwidth}{+2cm}
%\usepackage{showframe}

%\usepackage{lastpage}
% for \pageref{LastPage}

% fancy header
\usepackage{fancyhdr}
\pagestyle{fancy}
%\pagestyle{plain} % Imprime simplement le numero de page centre en bas
%\pagestyle{headings}
\fancyhf{} % clear default header and footer

\fancyhead[LE,RO]{\textit{\nouppercase{\leftmark}}}
%\fancyhead[LO,RE]{\textit{\nouppercase{\thepage{} of \pageref{LastPage}}}}
\fancyhead[CE,CO]{}
\fancyfoot[C]{--\thepage--}
\setlength{\headheight}{14pt}

\newcommand\Hrule{\noindent\rule{\linewidth}{1.5pt}}
%\AL: to add abs and norm for fracs
%\DeclarePairedDelimiter\abs{\lvert}{\rvert}%
%\DeclarePairedDelimiter\norm{\lVert}{\rVert}%
\usepackage[pdftex,
    pdfauthor={Alexandre Lafaye},
    pdftitle={Laboratoire PHY1 E8},
    pdfsubject={},
    pdfkeywords={rapport,laboratoire,HEIG-VD,PYH1,Thérmodynamique},
pdfcreator=pdflatex]{hyperref}
\hypersetup{colorlinks, citecolor=black, filecolor=black, linkcolor=black, urlcolor=black}

\newcommand{\weblink}[2]{\href{#1}{\url{#1}: \textsc{#2}}}

%Maths shortcuts
%\newcommand{\R}{\mathbb{R}}
%\newcommand{\E}{\mathbb{E}}
%\newcommand{\N}{\mathbb{N}}
%\newcommand{\Var}{\mathrm{Var}}
%\newcommand{\Cov}{\mathrm{Cov}}
%\newcommand{\Prob}{\mathbb{P}}
%\newcommand{\trace}{\mathbb{\text{trace}}}
%\def\mathbi#1{\textbf{\em #1}}
%%\def\mathi#1{\em #1}
%\newcommand{\Z}{\mathbi{Z}_2}
%\newcommand{\kMeans}{{\em k}-Means\xspace}
%\newcommand{\kmeans}{{\em k}-means\xspace}
%\newcommand{\kmeanspp}{{\em k}-means$++$\xspace}
%\newcommand{\captionwidth}{0.9\textwidth}

% index stuff
\usepackage{tocbibind} % Index in TOC
%\usepackage{makeidx}
%\usepackage{showidx} % show index entries in margin
%\makeindex

%\newcommand{\bib}[4]{\item{\textsc{#1}, \emph{#2}, #3\\\hspace{2em}#4}}

% ---------------------------------------------------------

\newcommand{\titleinfo}{Laboratoire de Physique E8 }
\author{Alexandre Lafaye}
\date{07 Avril 2016}
\title{\titleinfo}

% Margins for handwritten notes
%\addtolength{\oddsidemargin}{-2cm}
%\addtolength{\evensidemargin}{-2cm}
%\addtolength{\textwidth}{-2cm}

\usepackage{color}
\definecolor{Brown}{cmyk}{0,0.81,1,0.60}
\definecolor{OliveGreen}{cmyk}{0.64,0,0.95,0.40}
\definecolor{CadetBlue}{cmyk}{0.62,0.57,0.23,0}
\definecolor{lightlightgray}{cmyk}{0,0,0,0.02}
\definecolor{gray}{cmyk}{0,0,0,0.8}

% LaTeX magic: make sections and parts have a cleardoublepage
%\let\stdsection\section
%\renewcommand*\section{\cleardoublepage\stdsection}
\captionsetup{belowskip=12pt,aboveskip=4pt}
\begin{document}
% JP
%\frontmatter

\lstset{
    language=Python,                             % Code langugage
    basicstyle=\ttfamily,                   % Code font, Examples: \footnotesize, \ttfamily
    keywordstyle=\color{OliveGreen},        % Keywords font ('*' = uppercase)
    commentstyle=\color{gray},              % Comments font
    numbers=left,                           % Line nums position
    numberstyle=\tiny,                      % Line-numbers fonts
    stepnumber=1,                           % Step between two line-numbers
    numbersep=5pt,                          % How far are line-numbers from code
    backgroundcolor=\color{lightlightgray}, % Choose background color
    frame=none,                             % A frame around the code
    tabsize=2,                              % Default tab size
    captionpos=b,                           % Caption-position = bottom
    breaklines=true,                        % Automatic line breaking?
    breakatwhitespace=false,                % Automatic breaks only at whitespace?
    showspaces=false,                       % Dont make spaces visible
    showtabs=false,                         % Dont make tabls visible
    morekeywords={uint32_t,uint8_t,uint16_t,time_spec_t,size_t,
    clock_config_t},
}
\pagestyle{empty}
\pagenumbering{arabic}
\begin{titlepage}
    \null\vfil
    \begin{flushleft}
	\begin{center}
	    \Huge \textbf{\titleinfo}
	\end{center}
    \end{flushleft}
    \par
    \hrule height 4pt
    \par
    \begin{flushright}
	\large
	\textsl{Rapport de laboratoire} \par
    \end{flushright}
    \vspace{\fill}
    \vspace*{\stretch{1}}

    \begin{center}
	\Large
	\textsc{Colin Bairoch, Julia Cosandier, Alexandre Lafaye\\
	    %\normalsize
	HEIG-VD -- 2016} \\
	\href{mailto:julia.cosandier@heig-vd.ch}{julia.cosandier@heig-vd.ch}
	\href{mailto:colin.bairoch@heig-vd.ch}{colin.bairoch@heig-vd.ch}
	\href{mailto:alexandre.lafaye@epfl.ch}{alexandre.lafaye@epfl.ch}
    \end{center}
    \vspace{\fill}
    \vspace*{\stretch{2}}
    \begin{center}
	\large
	Supervisé par: \\
	\vspace{\fill}
	Ioan Balin -- HEIG-VD\\
	\href{http://heig-vd.ch/}{http://heig-vd.ch/} \\
    \end{center}
    \begin{center}
	\large
	Mirko Croci -- EPFL \\
	\href{http://heig-vd.ch/}{http://heig-vd.ch/} \\
    \end{center}

    \vspace{\fill}
    \vspace*{\stretch{2}}

    \begin{figure}[!h]
	\centering
	\parbox{1.4in}{
	\includegraphics[height=4em]{logo_heig_h80.jpg}}
    \end{figure}

    \vspace*{1cm}
\end{titlepage}

\pagestyle{fancy}

\tableofcontents

\newpage

\section{Introduction}
%\addcontentsline{toc}{section}{Introduction}
\markboth{Introduction}{}
\subsection*{Objectifs du Laboratoire}
Dans ce travail pratique les buts suivants ont été visés:
\begin{itemize}
    \item{Mesure du volume d'une colonne d'air en fonction de la pression (avec température constante)}
    \item{}
    \item{Confirmation de la loi de Boyle-Mariotte}
    \item{Tracés des régressions linéaires de la pression en fonction de l'inverse du volume à partir des mesures précédentes}
    \item{Détermination de }
\end{itemize}

\newpage

\section{Théorie}
%\addcontentsline{toc}{section}{Théorie}
\markboth{Théorie}{}
\subsection{Gaz parfaits}
Pour un gaz dit "parfait" présent en quantité n de moles, on peut définir son état à partir des 3 constantes suivantes:
\begin{itemize}
    \item{La pression p en Pascal}
    \item{Le volume V en mètre cube}
    \item{La température en degré Kelvin}
\end{itemize}

De manière plus formelle:
\begin{equation}
    p \cdot V = n\cdot R \cdot T [J]
\end{equation}

R est la constante des gaz: $R=8.314 [J\cdot K^{-1} \cdot mol^{-1}]$

Si l'une de ces grandeurs reste constante, les deux autres ne peuvent pas varier indépendamment l'une de l'autre. En fixant la température, la loi de Boyle-Mariotte est obtenue:
\begin{equation}
    \label{boylemariotte}
    p\cdot V = constante
\end{equation}

Cette loi peut être observée grâce à un thermomètre à gaz.
Le thermomètre à gaz est un instrument de mesure constitué d'un grand tube de verre scellé à l'une de ses extrémités. De l'autre côté, une goutte de mercure assure l'étanchéité. Ce mercure va donc se mouvoir en fonction du rapport entre la pression extérieure $p_0$ et intérieure et définir un volume intérieur $V_0$.\\
En pompant de l'air, il est possible de créer une dépression d'un côté du mercure et obtenir une pression: $p_0 \Delta p$.\\
Une pression est aussi exercée par l'air enfermé par le mercure:
\begin{equation}
    p_{Hg} = \rho_{Hg}\cdot g \cdot h_{Hg}
\end{equation}
Sachant que:
\begin{align*}
    &\rho_{Hg} = 13600 kg/m^3\\
    &g = 9.81 m/s^2\\
    &h_{Hg} = hauteur\ de\ la\ colonne\ de\ mercure
\end{align*}

Donc la pression dans la colonne d'air du côté scellé vaut:
\begin{equation}
    p = p_0 + p_{Hg} + \Delta p
\end{equation}

Pour trouver le volume associé, il suffit d'utiliser la hauteur et la section intérieure du tube de verre.
\begin{equation}
    V = \frac{\pi \cdot d^2 \cdot h}{4}
\end{equation}
Le diamètre intérieur du tube vaut: $d=0.27cm$.

Quelques rappels concernant les conversions d'unités de pression:
\begin{align*}
    &1Pa = 1 N/m^2\\
    &1bar = 10^5 Pa\\
    &1torr = 1mmHg = 133.3224 Pa\\
    &760mmHg = 760torr=1atm=101325Pa=1.01325bar
\end{align*}

\subsection{Régression linéaire}

Pour la regression linéaire, les formules suivantes peuvent être utilisées afin de calculer les incertitudes liées à la courbe générée.
Dans ces équations, n représente le nombre de mesures et p la pente.

\begin{equation}
    \label{regressionr}
    r=\frac{\sum_{k=1}^{n}(x_k-\bar{x})\cdot(y_k-\bar{y})}{\sqrt{\sum_{k=1}^{n}(x_k-\bar{x})^2\cdot\sum_{k=1}^{n}(y_k-\bar{y})^2}}
\end{equation}
\\

\begin{equation}
    \label{deltam}
    \Delta p=p\sqrt{\frac{\frac{1}{r^2}-1}{n-2}}
\end{equation}

\begin{equation}
    \label{deltah}
    \Delta h = \Delta p \cdot \sqrt{\frac{\sum_{k=1}^{n}x_k^2}{n}}
\end{equation}

\newpage

\section{Manipulation}
%\addcontentsline{toc}{section}{Manipulation}
\markboth{Manipulation}{}
\subsection{Préparation de la manipulation}
Cette expérience demande de préparer le thermomètre à gaz en quelques étapes simples.
\begin{enumerate}
    \item{La pompe à main doit être branchée avec le tube ayant son extremité scellée orientée vers le haut}
    \item{Tapoter le tube afin de réunir en une seule goutte les éclats de mercure disposés le long du tube}
    \item{Avec la pompe, générer une mesure de $\Delta p$ maximale}
\end{enumerate}

Pour régler le volume initial $V_0$:
\begin{enumerate}
    \item{Retourner le thermomètre à gaz pour que la partie scellée se retrouve en bas}
    \item{En controlant la pression avec la vanne du thermomètre à main, réduire la dépression $\Delta p$ jusqu'à 0}
\end{enumerate}

\subsection{Méthode}
Pour étudier la loi de Boyle-Mariotte, des mesures à trois températures différentes ont été effectuées:
\begin{itemize}
    \item{$T_1 = 280K = 6.85C\degree$}
    \item{$T_2 = T_{ambiante}$}
    \item{$T_3 = 340K = 66.85C\degree$}
\end{itemize}

Pour les cas 280K et 340K l'eau à été amenée à la température voulue dans un cas en la mélangeant avec des glaçons et en attendant que la température remonte, dans l'autre en faisant bouillir l'eau et en la mélangeant avec du liquide à température ambiante.

Pour chaque série de mesures, la température a été mesurée avant et après.

La procédure a donc été la suivante pour chaque température:
\begin{itemize}
    \item{Le thermomètre à gaz a été mis à $\Delta p =0$}
    \item{La hauteur de colonne de mercure a été mesurée}
    \item{Avec la pompe à main une dépression $\Delta p$ a été générée aux valeurs se trouvant dans les tableaux}
    \item{La hauteur de la colonne d'air a été mesurée, après qu'elle se soit stabilisée}
\end{itemize}

\subsection{Résultats}

Pour tout ce qui va suivre, le baromètre du laboratoire a donné comme pression atmosphérique le jour de la manipulation:
\begin{align}
    &P_{atm} = 969mbar
    &P_{atm} = 969.6mbar
    &P_{atm} = 969.8mbar
\end{align}

\paragraph{Temperature T=280K}
Voici donc les résultats bruts pour une température visée de $6.85C\degree$. Toutefois, selon les mesures de la température initiale et finales:

\begin{align*}
    &T_{ini} = 7.5C\degree\\
    &T_{fin} = 11C\degree
\end{align*}

On peut donc affirmer qu'en moyenne $T_{moy} = 9.25C\degree$.

Toutefois, en première approximation, on peut dire que la chaleur s'est diffusé de manière constante, et donc que la température a augmenté de manière linéaire.
En échelonnant le abcisses avec le nombre de mesures, on obtient une droite de valeur:
\begin{equation}
    \frac{3}{8}\cdot x + 7.5
\end{equation}

Il suffit ensuite de remplacer x par le numéro de la mesure afin d'obtenir une température plus proche de la réalité au moment de la températeure.
Le défaut de ce modèle est qu'il suppose que la diffusion de chaleur est une fonction linéaire (ce qui généralement est faux) et que les mesures ont été prisese de manière régulière dans le temps (ce qui est aussi inexact). Mais cela reste une meilleure approximation que la moyenne.

\begin{table}[h]
    \centering
    \caption{Résultats T=280K}
    \begin{tabular}{|l|l|l|}
	\hline
	$\Delta p [Pa]$	&Hauteur [m] &Temperature [C\degree]\\
	\hline
	100	&10.3 & 7.5   \\
	200	&11.5 & 7.88  \\
	300	&13.1 & 8.25  \\
	400	&15.3 & 8.63  \\
	500	&17.7 & 9     \\
	600	&21.1 & 9.38  \\
	650	&23.7 & 9.75  \\
	700	&26.6 & 10.13 \\
	720	&28.2 & 10.50 \\
	740	&29.2 & 10.88 \\
	760	&32.1 & 11.25 \\
	780	&33.6 & 11.63 \\
	\hline
    \end{tabular}
\end{table}

\paragraph{Température ambiante}
\begin{table}[h]
    \centering
    \caption{Résultats $T=T_{ambiante}$}
    \begin{tabular}{|l|l|}
	\hline
	$\Delta p [Pa]$	&Hauteur [m]\\
	\hline
	100	&11.4\\
	200	&12  \\
	300	&13.7\\
	400	&15.6\\
	500	&18.4\\
	600	&22.1\\
	650	&25  \\
	700	&27.6\\
	720	&29.1\\
	740	&30.5\\
	760	&32.9\\
	780	&34.9\\
	\hline
    \end{tabular}
\end{table}

\paragraph{Temperature T=340K}
Voici donc les résultats bruts pour une température visée de $66.85C\degree$. Toutefois, selon les mesures de la température initiale et finales:

\begin{align*}
    &T_{ini} = 66C\degree\\
    &T_{fin} = 50C\degree
\end{align*}

On peut donc affirmer qu'en moyenne $T_{moy} = 58C\degree$.

Cette fois ci, l'approximation linéaire est d'autant plus importante que le delta entre la température initiale et finale est grande. Cette courbe est décroissante:
\begin{equation}
    -\frac{4}{3}\cdot x + 66
\end{equation}

\begin{table}[h]
    \centering
    \caption{Résultats T=340K}
    \begin{tabular}{|l|l|l|}
	\hline
	$\Delta p [Pa]$	&Hauteur [m] &Temperature[C\degree]\\
	\hline
	100	&12.2 &   66   \\
	200	&13.4 &   64.67\\
	300	&15.5 &   63.33\\
	400	&17.8 &   62   \\
	500	&20.7 &   60.67\\
	600	&24.8 &   59.33\\
	650	&27.5 &   58   \\
	700	&31.3 &   56.67\\
	720	&32.5 &   55.33\\
	740	&33   &   54   \\
	760	&\~34   & 52.67\\
	780	&\~34.1 & 51.33\\
	\hline
    \end{tabular}
\end{table}

\newpage

\section{Conclusion}
%\addcontentsline{toc}{section}{Conclusion}
\markboth{Conclusion}{}
\input{concl}

\section{Annexe}
%\addcontentsline{toc}{section}{Annexe}
\markboth{Annexe}{}
\input{annexe}

\listoffigures
\listoftables

\end{document}
